% ═══════════════════════════════════════════════════════════════
% QR-CODE: Minitest Pretérito Imperfecto
% Spanisch Klasse 12
% ═══════════════════════════════════════════════════════════════

\documentclass[11pt, a4paper]{article}

\usepackage[utf8]{inputenc}
\usepackage[T1]{fontenc}
\usepackage[ngerman]{babel}
\usepackage{helvet}
\renewcommand{\familydefault}{\sfdefault}

\usepackage[margin=2cm]{geometry}
\usepackage{xcolor}
\usepackage{qrcode}
\usepackage{tikz}
\usepackage[hidelinks]{hyperref}

\definecolor{akzent}{HTML}{C62828}

% URL als Variable definieren
\newcommand{\mturl}{https://devbox42.github.io/sciencesim/spanisch/kl12/preterito-imperfecto/MT-01-preterito-imperfecto.html}

\pagestyle{empty}

\begin{document}

\begin{center}

\vspace*{2cm}

{\Huge\bfseries\textcolor{akzent}{Minitest}}

\vspace{0.5cm}

{\LARGE Pretérito Imperfecto}

\vspace{0.3cm}

{\large Spanisch Klasse 12 | 30 BE}

\vspace{2cm}

% QR-Code (klickbar im PDF)
\href{\mturl}{\qrcode[height=6cm]{\mturl}}

\vspace{1.5cm}

{\Large Scanne den QR-Code mit deinem Tablet}

\vspace{0.5cm}

{\small oder öffne den Link im Browser}

\vspace{2cm}

\begin{tikzpicture}
\draw[akzent, line width=2pt] (0,0) -- (10,0);
\end{tikzpicture}

\vspace{0.5cm}

{\footnotesize\textcolor{gray}{MT-01-preterito-imperfecto | Spanisch Klasse 12}}

\end{center}

\end{document}
